\begin{table}[!htbp]\centering
\caption{Income's contribution conditional on scores, year, and home region}
\label{tab:region-d5-high-math-main-cubic-income-contribution}
\small
\begin{tabular}{lrrrr}\hline\hline
Outcome & $R^2$ without income & Full $R^2$ & Income $\Delta R^2$ & Income share (\%) \\ \hline
Premium institution & 0.3101 & 0.3297 & 0.0196 & 5.94 \\
Premium field & 0.0746 & 0.0870 & 0.0124 & 14.26 \\
Either (or both) & 0.1404 & 0.1411 & 0.0007 & 0.48 \\
Both & 0.1736 & 0.1741 & 0.0005 & 0.30 \\
\hline\hline\end{tabular}
\begin{minipage}{0.98\textwidth}\footnotesize
\textit{Notes:} Income share is the increase in ordinary unadjusted $R^2$ from adding nine income-decile indicators, divided by full-model $R^2$. Both models retain all nine total-degree-three math/verbal terms and admission-year indicators, using the same 93,423 students. Selection requires original $z_{math} \geq 1$. Both models include reported home-region indicators; the full model uses income decile 5 as reference. This is the unique incremental contribution after scores, admission year, and home region; shared explanatory variation is not allocated to income. The 840 admission-2026 records remain. All fit measures are in-sample and descriptive.
\end{minipage}
\end{table}
